% Terminology and editorial conventions
% ========================================
%
% The policy is the repository's:
%
%   1. Use standard modern English where it exists, so that a notion can be
%      found in the current literature under the word used here.
%   2. Where the French term is a coinage or has no settled English
%      counterpart, give an accurate rendering and record it below.
%   3. Translate sentences, not paragraphs: definitions keep the original's
%      shape (data, then axioms) so that the translation can be checked
%      against the French clause by clause.
%
% The authors' own English abstract uses "n-stack", "prestack", "descent
% data", "closed model category (cmc)" and "left Quillen presheaf"; those
% spellings are authoritative and are used throughout.
%
% Bibliographic titles are left in the language in which they were printed;
% only the authors' annotations in the reference list are translated.  A
% translated title would not be searchable, and inventing one would be a
% silent editorial intervention.

\begin{table}[htbp]
  \centering
  \small
  \begin{tabular}{p{0.36\textwidth} p{0.30\textwidth} p{0.24\textwidth}}
    \hline
    \textbf{French (original)} & \textbf{English (this translation)} & \textbf{Note}\\[2pt]
    \hline
    $n$-champ                     & $n$-stack                      & \\
    préchamp                      & prestack                       & \\
    champ                         & stack                          & \\
    champ de modules              & moduli stack                   & \\
    champ associé                 & associated stack               & \\
    $n$-champ de Segal            & Segal $n$-stack                & \\
    $n$-préchamp de Segal         & Segal $n$-prestack             & \\
    $n$-précat de Segal           & Segal $n$-precategory          & see note below\\
    catégorie de Segal            & Segal category                 & \\
    catégorie simpliciale         & simplicial category            & \\
    ensemble simplicial           & simplicial set                 & \\
    catégorie de modèles fermée   & closed model category (cmc)    & authors' abbreviation\\
    cmf                           & cmc                            & \\
    structure de cmf              & cmc structure                  & \\
    catégorie de modèles interne  & internal model category        & \\
    donnée de descente            & descent data                   & \\
    donnée de descente généralisée & generalised descent data (gdd) & \\
    effectif / effectivité        & effective / effectivity        & \\
    recollement / recoller        & gluing / to glue               & \\
    crible                        & sieve                          & \\
    topologie grossière           & coarse topology                & \\
    site                          & site                           & \\
    site annelé                   & ringed site                    & \\
    préfaisceau                   & presheaf                       & \\
    faisceau                      & sheaf                          & \\
    hyper-recouvrement            & hypercovering                  & \\
    préfaisceau de Quillen        & Quillen presheaf               & \\
    \quad à gauche                & \quad left                     & \\
    foncteur de Quillen           & Quillen functor                & \\
    équivalence faible            & weak equivalence                & \\
    équivalence objet-par-objet   & object-by-object equivalence   & \\
    remplacement fibrant          & fibrant replacement            & \\
    objet fibrant / cofibrant     & fibrant / cofibrant object     & \\
    localisation de Dwyer--Kan    & Dwyer--Kan localisation        & \\
    strictification               & strictification                & \\
    troncation                    & truncation                     & \\
    intérieur                     & interior                       & \\
    section équilibrée (eq-section) & balanced section (eq-section) & authors' term\\
    fibré universel               & universal bundle               & \\
    famille universelle           & universal family               & \\
    complexe parfait              & perfect complex                & \\
    support cohomologique         & cohomological support          & \\
    \hline
  \end{tabular}
  \caption{Terminology.  Page references, when needed, are to the original
  PDF of arXiv:math/9807049v3 (1--250).}
  \label{tab:terminology}
\end{table}

% NAMES AND NOTATION KEPT UNCHANGED
% --------------------------------
% "n-précat" is the authors' abbreviation of "n-précatégorie", not of
% "n-préchamp".  The two are distinct: nSePC is the category of Segal
% n-precategories (simplicial objects with the constancy conditions but not
% necessarily the Segal conditions), whereas a Segal n-prestack is a presheaf
% with values in nSePC.  Translating both as "prestack" collapses the
% distinction; the translation keeps "precategory" for the first and
% "prestack" for the second throughout.
%
% "section équilibrée", abbreviated "eq-section", is the authors' own term
% (introduced in the Introduction, defined in §16).  The "eq" abbreviates
% French "équilibrée" ("balanced"), not "equi-free" or "equivariant"; there
% is no established English term for it.  We use "balanced section
% (eq-section)" at the single occurrence in the Introduction, with a
% translator's note, and keep "eq-section" in §16 and later.
%
% Fl(A)          the Segal (n-1)-category of 1-morphisms of A
% Fl^i(A)(x,y)   the subcategory of i-morphisms with source x and target y
% nSeCAT, nSeCHAMP, nCHAMP, nSeCAT', Cpx_O, qis
% \Xx, \Zz, \Gg, \Oo, \Bb, \mm  -- the authors' own macros
%
% "cmc" is kept rather than "model category" because the authors use it
% systematically, including in their English abstract, and because the
% phrase "closed model category" is part of Quillen's terminology.
%
% "0-champ de Segal" is rendered "Segal 0-stack", not "0-Segal stack":
% the adjective modifies the whole notion, as in "Segal n-category".
